\documentclass[11pt,a4paper]{article}
\usepackage[utf8]{inputenc}
\usepackage{amsmath}
\usepackage{amsfonts}
\usepackage{hyperref}
\usepackage{geometry}
\geometry{margin=1in}

\title{Technical Note: Parameter Centering around a Typical Operating Point}
\author{FSA-v5 Research Engineering}
\date{May 10, 2026}

\begin{document}

\maketitle

\section{Introduction}
In the FSA-v5 model, the non-linear SDE dynamics are defined using constants $A_{TYP}$ (Typical Alertness) and $F_{TYP}$ (Typical Fatigue). This technique is known in statistics and control theory as \textbf{Parameter Centering} or \textbf{Operating Point Analysis}. This note explains why this technique is used and its impact on model identifiability.

\section{The Mathematical Mechanism}

\subsection{Quadratic Fatigue Centering}
The Stuart-Landau bifurcation parameter $\mu$ includes a quadratic penalty on fatigue. Without centering, the term is written as $\mu \propto - \mu_F F - \mu_{FF} F^2$. This form creates high correlation between the linear and quadratic coefficients, making Bayesian inference (SMC$^2$) unstable.

In FSA-v5, we center the term on the typical operating point $F_{TYP} = 0.20$:
\begin{equation}
    \mu(F) \propto - \mu_{FF} (F - F_{TYP})^2
\end{equation}
This centering decorrelates the parameters. $\mu_{FF}$ now isolatedly controls the \textbf{curvature} of the "fatigue cliff" specifically at the intensity levels where the athlete actually operates.

\subsection{Autonomic Boost Normalization}
The model assumes that adaptation gains are modulated by the current autonomic state $A$. To ensure that the gains represent "deviations from normal," we normalize the boost factor by the typical operating point $A_{TYP} = 0.10$:
\begin{equation}
    a_{B}(A) = \frac{1 + \epsilon_{AB} A}{1 + \epsilon_{AB} A_{TYP}}
\end{equation}
By construction, when $A = A_{TYP}$, the multiplier $a_B = 1.0$, and the system behaves according to standard training expectations.

\section{Advantages of Centering}

\begin{itemize}
    \item \textbf{Identifiability:} By subtracting the typical value, we reduce the multi-collinearity of the parameter space. This allows the filter to uniquely identify fitness gains vs. fatigue sensitivity from noisy data.
    \item \textbf{Numerical Stability:} Centering avoids "catastrophic cancellation" in floating-point arithmetic, which occurs when subtracting two nearly-equal large numbers to find a small difference.
    \item \textbf{Intuitive Prior Biasing:} Prior distributions for parameters like $\mu_0$ (baseline drive) become meaningful; they describe the state at a \textbf{healthy human baseline} rather than at a hypothetical zero-state (which is physiologically impossible).
\end{itemize}

\section{Conclusion}
The use of $A_{TYP}$ and $F_{TYP}$ is not a specific model variant, but a robust engineering strategy to align abstract non-linear dynamics with the actual operating regime of human physiology. It transforms a mathematically difficult optimization problem into a stable, identifiable one.

\end{document}
